The spectra of two stars with different temperatures $T_1$ and respectively
$T_2$ were compared. In the spectrum of each star, two very close spectral
lines corresponding to the wavelength with values $\lambda_1$ and respectively
$\lambda_2$ were found. For each line of this spectral lines, the difference
between the corresponding visual apparent magnitudes of the stars are
$\Delta m_{\lambda_1} = m_{1,\lambda_1} - m_{2,\lambda_1}$ and
$\Delta m_{\lambda_2} = m_{1,\lambda_2} - m_{2,\lambda_2}$.
$m_{1,\lambda_1}$ is the apparent magnitude of the star 1 for the wavelength
$\lambda_1$, $m_{1,\lambda_2}$ is the apparent magnitude of the star 1 for the
wavelength $\lambda_2$, $m_{2,\lambda_1}$ is the apparent magnitude of the
star 2 for the wavelength $\lambda_1$, $m_{2,\lambda_2}$ is the apparent
magnitude of the star 2 for the wavelength $\lambda_2$.

\emph{Determine} the temperature $T_1$ of one of the two stars, if the
temperature $T_2$ of the other star is already known, by using the Planck
expression of black body radiation:
\[ r(\lambda) = \frac{2\pi h c^2}{\lambda^5}
   \left( e^{hc/k\lambda c} - 1 \right)^{-1}, \]
% sic: the original paper prints the exponent as hc/k\lambda c (presumably a
% typo for hc/k\lambda T; cf. the closing condition below).
where: $h$ --- Planck's constant; $k$ --- Boltzmann's constant; $c$ --- speed
of light in vacuum. You will consider that $hc \gg k\lambda T$.
